
\documentclass[12pt,a4paper]{article}
\usepackage[utf8]{inputenc}
\usepackage[T1]{fontenc}
\usepackage{geometry}
\geometry{margin=1in}
\usepackage{helvet}
\usepackage{xcolor}
\usepackage{titlesec}
\usepackage{booktabs}
\usepackage{hyperref}

% Define colors
\definecolor{darkblue}{RGB}{0, 51, 102}

% Formatting headers
\titleformat{\section}{\large\bfseries\color{darkblue}}{}{0em}{}[\titlerule]
\titleformat{\subsection}{\bfseries\color{darkblue}}{}{0em}{}

\begin{document}

\begin{center}
    {\LARGE \textbf{Final Project Proposal}} \\
    \vspace{0.5cm}
    {\large \textbf{CSE 429: Computer Vision and Pattern Recognition}} \\
    \vspace{1cm}
    \textbf{Project Title: Benchmarking DINO for Drone-View Object Detection on VisDrone} \\
    \vspace{1cm}
    \textbf{Submitted to:} Prof. Ahmed Gomaa \\
    \vspace{0.5cm}
    \textbf{Date:} \today
\end{center}

\vspace{1cm}

\section{Team Members}
\begin{itemize}
    \item Jana Mohamed Elsalhy --- 120220093
    \item Shiref Ashraf Mohamed --- 120220099
    \item Ahmed Nagah Ramadan --- 120220116
    \item Tasbih Othman Neamatallah --- 120220117
    \item Karim Yaser Mazroua --- 120220118
    \item Yasmeen Sameh Mohamed --- 120220143
\end{itemize}

\section{Problem Statement}
Aerial object detection from drone-mounted cameras is a high-impact task characterized by unique environmental constraints. The system must process raw aerial frames as input and output precise bounding boxes, class labels, and confidence scores for diverse targets (e.g., pedestrians, vehicles, bicycles).

\textbf{The Challenge:} 
Standard detection models often fail in UAV contexts because objects appear at \textbf{extremely small scales} (often $< 10 \times 10$ pixels), perspective shifts are extreme, and scene density is high. Reliable detection is critical for military surveillance, search and rescue, and traffic monitoring. This project evaluates whether the \textbf{DINO} (DETR with Improved DeNoising) architecture can bridge the gap between high-precision transformer logic and the specific demands of drone-view imagery.

\section{Approach}
We adopt the "Review and Implement" track, comparing three distinct architectures on the \textbf{VisDrone} dataset:
\begin{enumerate}
    \item \textbf{DINO (Primary):} An anchor-free, end-to-end transformer detector utilizing Contrastive DeNoising (CDN). We will fine-tune the official model using pretrained COCO checkpoints.
    \item \textbf{Vanilla DINO:} A baseline version without the denoising components to isolate the impact of the denoising mechanism (Ablation Study).
    \item \textbf{YOLOv11:} A CNN-based real-time detector to benchmark against the current industry standard for edge-deployment.
\end{enumerate}

\section{Experiments and Results}
\subsection{Dataset}
We will use the \textbf{VisDrone-DET2019} dataset, containing over 10,000 images and 10 object categories collected from drone cameras.

\subsection{Implementation Details}
\begin{itemize}
    \item \textbf{Framework:} PyTorch, utilizing the official DINO and Ultralytics YOLOv11 repositories.
    \item \textbf{Training:} Due to compute constraints, we will perform prototype-level training for ~20 epochs.
    \item \textbf{Hardware:} NVIDIA T4 (Kaggle/Colab).
\end{itemize}

\subsection{Planned Experiments}
\begin{table}[h]
\centering
\begin{tabular}{lp{9cm}}
\toprule
\textbf{Experiment} & \textbf{Purpose} \\ \midrule
DINO Fine-tuning & Primary performance results on drone imagery. \\
Vanilla Ablation & Quantify the impact of the denoising mechanism. \\
YOLOv11 Benchmark & Establish a baseline for real-time speed/efficiency. \\ \bottomrule
\end{tabular}
\end{table}

\section{Success Criteria}
Success is defined by the functional implementation of the DINO pipeline on VisDrone, achieving a higher mAP than the vanilla baseline, and providing a comprehensive trade-off analysis between DINO's accuracy and YOLOv11's speed for real-world drone applications.

\end{document}
